\chapter{Preface}

These notes build a mathematical path toward Tilo's roadmap on clustered
excitatory/inhibitory spiking networks, mesoscopic theory, multistability, and
chaos. The goal is not to summarize every model in isolation. The goal is to
teach the sequence of ideas that makes the central research question precise:
when a clustered spiking network switches among population states, is the
switching a finite-size escape among stable attractors, or is it deterministic
mesoscopic chaos generated by frozen inter-cluster structure?

The course begins with dynamical systems because attractors, stability,
bifurcations, and Lyapunov exponents are the language of the distinction. It
then adds spike trains, point processes, filtering, and finite-size scaling so
that the stochastic objects in a spiking network can be treated honestly. The
middle chapters introduce single-neuron models, balanced E/I recurrence,
clustered attractor networks, random rate chaos, and dynamic mean-field theory.
The final chapters specialize to paired E/I clustered networks and develop the
diagnostics needed to separate multistability from chaos in simulations.

The source roadmap PDF is treated as a research plan, not as a polished
textbook. Ambiguous source claims are marked in the coverage audit under
\texttt{source/roadmap-coverage.md} and \texttt{source/pdf-extraction-notes.md}.
In particular, the suspicious PDF reference [4] is retained in the source
BibTeX for traceability but is not used in these notes as evidence for neural
chaos. All figures in this course are original teaching schematics generated
for these notes; none are copied or redrawn from the source PDF.
